% =============================================================================
% Applications of Algorithms (COMS4032A / COMS7041A) -- Assignment 2
% Binary Search Trees and Order-Statistic Trees
% Build:  pdflatex report.tex && pdflatex report.tex   (figures in ./figs/)
% =============================================================================

\documentclass[conference]{IEEEtran}
\IEEEoverridecommandlockouts

\usepackage{cite}
\usepackage{amsmath,amssymb}
\usepackage{algorithm}
\usepackage{algpseudocode}
\usepackage{graphicx}
\usepackage{booktabs}
\usepackage{siunitx}
\usepackage{xcolor}
\usepackage[hidelinks]{hyperref}

\graphicspath{{figs/}}

\newcommand{\INS}{\textsc{Tree-Insert}}
\newcommand{\DEL}{\textsc{Tree-Delete}}
\newcommand{\WALK}{\textsc{Inorder-Tree-Walk}}
\newcommand{\SEL}{\textsc{OS-Select}}
\newcommand{\RANK}{\textsc{OS-Rank}}
\newcommand{\TRANS}{\textsc{Transplant}}
\newcommand{\NIL}{\textsc{nil}}
\newcommand{\new}{\textcolor{blue!60!black}{$\blacktriangleleft$ new}}
\algrenewcommand\algorithmiccomment[1]{\hfill{\footnotesize #1}}

\begin{document}

\title{Empirical Analysis of Randomly Built Binary Search Trees and Order-Statistic Trees}

\author{\IEEEauthorblockN{Fortune Mnisi}
\IEEEauthorblockA{Student number: 1903697\\
School of Computer Science and Applied Mathematics\\
University of the Witwatersrand, Johannesburg\\
COMS4032A/COMS7041A: Applications of Algorithms}}

\maketitle

\begin{abstract}
The binary search tree operations \INS{}, \DEL{} and \WALK{} from Cormen
\textit{et al.}~\cite{clrs}, and an order-statistic tree built on a plain
(unbalanced) binary search tree supporting \SEL{} and \RANK{}, were implemented
in Java and evaluated on 18\,612 randomly built trees with up to
$n = 2^{22}$ keys. The mean tree height grew as $2.86\lg n - 7.7$, consistent
with Theorem~12.4 and with the slope of the known asymptotic expansion over the
same range, and stayed below the bound derived in the textbook at every size.
Instrumented runs showed that building and destroying a tree visit
$\Theta(n\lg n)$ nodes, matching their exact expected counts to within about 3\%.
Running times, however, grew faster than the operation counts: the cost of
visiting one node rose from about \SI{6}{\nano\second} to
\SI{100}{\nano\second} once the tree outgrew the processor caches. The same
effect made the $\Theta(n)$ inorder walk scale as $n^{1.41}$; laying the nodes
out in key order cut this to $n^{1.23}$ and gave a $5.3\times$ speedup.
Maintaining subtree sizes added almost nothing to insertion but made deletion
of a random node about twice as slow on large trees, because plain \DEL{}
touches only 1.75 nodes on average while the augmented version must update
every ancestor up to the root. All four order-statistic operations did
$\Theta(\lg n)$ work, with their time again dominated by cache misses.
\end{abstract}

\begin{IEEEkeywords}
binary search trees, order-statistic trees, randomly built trees, tree height,
empirical algorithmics, memory locality
\end{IEEEkeywords}

% -----------------------------------------------------------------------------
\section{Introduction}

A binary search tree (BST) supports its basic dynamic-set operations in $O(h)$
time, where $h$ is the height of the tree. The height depends entirely on the
order in which keys are inserted: sorted input produces a chain with
$h = n-1$, while a random insertion order produces, on average, a much shallower
tree. Theorem~12.4 (Section~12.4 of~\cite{clrs}) makes this precise, stating that a randomly
built BST on $n$ distinct keys has expected height $O(\lg n)$.

This report tests that claim empirically, measures the cost of building,
destroying and walking random BSTs, and then augments the BST with a
\emph{size} attribute to support order statistics. Throughout, the number of
nodes an operation visits is separated from the time it takes, since on real
hardware the two do not always grow at the same rate.

% -----------------------------------------------------------------------------
\section{Algorithms and Theoretical Analysis}

\subsection{Binary Search Tree Operations}

\WALK{} (Section~12.1 of~\cite{clrs}) and \INS{} and \DEL{}, with its
helper \TRANS{} (Section~12.3), were coded exactly as given in the textbook,
with \texttt{null} playing the role of
\NIL{}. \INS{} follows a single root-to-leaf path and takes $O(h)$ time. Given a
pointer to the node $z$, \DEL{} does constant work apart from the search for
$z$'s successor when $z$ has two children, so it takes $O(h)$ time in the
worst case but only time proportional to the height of $z$'s right subtree in
general. \WALK{} visits each node once and runs in $\Theta(n)$ time
regardless of the tree's shape (Theorem~12.1 of~\cite{clrs}).

\subsection{Height of a Randomly Built BST}

A randomly built BST is one produced by inserting $n$ distinct keys in an
order chosen uniformly from the $n!$ permutations. Theorem~12.4 bounds the
expected height via the exponential height $Y_n = 2^{X_n}$, showing that
\begin{equation}
E[Y_n] \le \tfrac{1}{4}\tbinom{n+3}{3}
\;\Rightarrow\;
E[X_n] \le \lg\!\left(\tfrac{1}{4}\tbinom{n+3}{3}\right),
\label{eq:clrs}
\end{equation}
which is about $3\lg n$. Sharper results are known. Devroye~\cite{devroye} showed
$E[X_n] \sim \alpha \ln n$ with $\alpha \approx 4.311$, and Reed~\cite{reed}
refined this to
\begin{equation}
E[X_n] = \alpha\ln n - \beta\ln\ln n + O(1),
\quad \beta = \frac{3}{2\ln(\alpha/2)} \approx 1.953.
\label{eq:reed}
\end{equation}
Since $\alpha\ln n \approx 2.988\lg n$, the textbook bound~\eqref{eq:clrs} is
close to tight in its leading constant. The $-\beta\ln\ln n$ term lowers the
local slope with respect to $\lg n$ to
$\ln 2\,(\alpha - \beta/\ln n)$, which is $2.87$ at $n = 2^{16}$.

\subsection{Building, Destroying and Walking}

Building a tree with $n$ calls to \INS{} visits, for each key, every ancestor
of the position where it lands. The total number of nodes visited is
therefore the internal path length of the final tree, whose expectation for a
random BST is
\begin{equation}
E[\text{build visits}] = 2(n+1)H_n - 4n \sim 2n\ln n \approx 1.386\,n\lg n,
\label{eq:build}
\end{equation}
where $H_n$ is the $n$th harmonic number~\cite{knuth}.

Destroying the tree by calling \DEL{} on the root until it is empty costs, per
call, a visit to the root plus, when the root has two children, a walk down
the left spine of its right subtree to find the successor. Deleting with the
successor in this way is known to leave the remaining tree random
(Theorem~H in Section~6.2.2 of~\cite{knuth}, after Hibbard~\cite{hibbard}).
When the root of a random tree of size $m$ has rank $r$, its right subtree is a
random tree of size $m-r$, whose left spine has expected length $H_{m-r}$.
Averaging over $r$ and summing over $m$ gives
\begin{equation}
E[\text{destroy visits}] = \sum_{m=1}^{n}\Bigl(1 + \frac{1}{m}\sum_{k=1}^{m-2}H_k\Bigr) \sim n\ln n,
\label{eq:destroy}
\end{equation}
about half the cost of building. Both operations should therefore take
$\Theta(n\lg n)$ expected time, with destruction roughly twice as fast.

\subsection{Order-Statistic Trees}

Following Section~14.1 of~\cite{clrs}, where \SEL{} and \RANK{} are given, each node $x$ stores
$x.\mathit{size}$, the number of nodes in the subtree rooted at $x$, so that
$x.\mathit{size} = x.\mathit{left}.\mathit{size} + x.\mathit{right}.\mathit{size} + 1$.
The textbook builds its order-statistic tree on a red-black tree. Here the
underlying tree is a plain BST, which changes three things.

\subsubsection*{Missing sentinel}
In a red-black tree every \NIL{} child is the sentinel $T.\mathit{nil}$, whose
size is 0. A plain BST has no sentinel, so every read of
$x.\mathit{left}.\mathit{size}$ or $x.\mathit{right}.\mathit{size}$ is replaced
by $\textsc{Size}(x.\mathit{left})$, which returns 0 for \NIL{}. This is the
only change needed in \SEL{} and \RANK{}, apart from \SEL{} returning \NIL{} if
it is asked for a rank outside $[1, n]$ instead of reading a field of \NIL{}.

\subsubsection*{Maintaining size during insertion}
\INS{} places the new node $z$ as a leaf. The nodes whose subtrees gain $z$
are exactly the nodes on the path from the root to $z$'s parent, which are the
nodes the search loop already visits. Incrementing each one as the loop passes
it, and giving $z$ size 1, keeps every size correct
(Algorithm~\ref{alg:ins}). The extra cost is one increment per node visited,
so the running time stays $O(h)$.

\begin{algorithm}[t]
\caption{\INS{}$(T, z)$ with size maintenance}
\label{alg:ins}
\begin{algorithmic}[1]
\State $y = \NIL$
\State $x = T.\mathit{root}$
\While{$x \ne \NIL$}
  \State $x.\mathit{size} = x.\mathit{size} + 1$ \Comment{\new}
  \State $y = x$
  \If{$z.\mathit{key} < x.\mathit{key}$} $x = x.\mathit{left}$
  \Else{} $x = x.\mathit{right}$
  \EndIf
\EndWhile
\State $z.\mathit{size} = 1$ \Comment{\new}
\State $z.p = y$
\If{$y == \NIL$} $T.\mathit{root} = z$
\ElsIf{$z.\mathit{key} < y.\mathit{key}$} $y.\mathit{left} = z$
\Else{} $y.\mathit{right} = z$
\EndIf
\end{algorithmic}
\end{algorithm}

\begin{algorithm}[t]
\caption{\DEL{}$(T, z)$ with size maintenance}
\label{alg:del}
\begin{algorithmic}[1]
\If{$z.\mathit{left} == \NIL$}
  \State \textsc{Decrement-To-Root}$(z.p)$ \Comment{\new}
  \State \TRANS$(T, z, z.\mathit{right})$
\ElsIf{$z.\mathit{right} == \NIL$}
  \State \textsc{Decrement-To-Root}$(z.p)$ \Comment{\new}
  \State \TRANS$(T, z, z.\mathit{left})$
\Else
  \State $y = \textsc{Tree-Minimum}(z.\mathit{right})$
  \State \textsc{Decrement-To-Root}$(y.p)$ \Comment{\new}
  \If{$y.p \ne z$}
    \State \TRANS$(T, y, y.\mathit{right})$
    \State $y.\mathit{right} = z.\mathit{right}$
    \State $y.\mathit{right}.p = y$
  \EndIf
  \State \TRANS$(T, z, y)$
  \State $y.\mathit{left} = z.\mathit{left}$
  \State $y.\mathit{left}.p = y$
  \State $y.\mathit{size} = z.\mathit{size}$ \Comment{\new}
\EndIf
\Statex
\Statex \textsc{Decrement-To-Root}$(x)$ \Comment{\new}
\While{$x \ne \NIL$}
  \State $x.\mathit{size} = x.\mathit{size} - 1$
  \State $x = x.p$
\EndWhile
\end{algorithmic}
\end{algorithm}

\subsubsection*{Maintaining size during deletion}
\DEL{} always removes exactly one node from one position in the tree: $z$
itself when $z$ has at most one child, or $z$'s successor $y$ when $z$ has two
children (in that case $y$ leaves its old position and takes $z$'s place). The
nodes that lose a descendant are precisely the ancestors of the position that
is vacated. Algorithm~\ref{alg:del} therefore walks from the parent of that
position up to the root, subtracting one from each size, before any pointers
are changed, while the parent pointers still describe the old path:
\begin{itemize}
\item If $z$ has at most one child, its child (possibly \NIL{}) moves up into
$z$'s place with its subtree unchanged. Only $z$'s ancestors lose a node, so
the walk starts at $z.p$.
\item If $z$ has two children, $y$ is removed from its old position, so the
walk starts at $y.p$. Because $y$ lies in $z$'s right subtree, this walk
passes through $z$ and then $z$'s ancestors (when $y.p = z$ it starts at $z$
itself). $y.\mathit{right}$ moves up into $y$'s old place with its subtree
unchanged. Finally $y$ takes over $z$'s position and $z$'s subtree, which now
holds every node $z$ had except $y$; that number is exactly $z$'s
already-reduced size, so $y.\mathit{size} = z.\mathit{size}$.
\end{itemize}
No other node's subtree changes, so all sizes stay correct. The upward walk
adds $O(h)$ work. This keeps \DEL{} within $O(h)$, but it changes its
character: the plain version only ever moves \emph{down} from $z$, while the
augmented version always walks the full distance from the vacated position up
to the root.

\subsection{Expected Cost of the Order-Statistic Operations}

\SEL{} follows one path down from the root to the node it returns, and
\RANK{} follows one path up from a node to the root, so both run in $O(h)$.
The augmented \INS{} and \DEL{} are also $O(h)$. On a random BST every one of
these paths has the length of a typical root-to-node path, whose expected
number of nodes is
\begin{equation}
E[\text{depth}] + 1 = 2\left(1 + \tfrac{1}{n}\right)H_n - 3 \sim 1.386\lg n,
\label{eq:depth}
\end{equation}
so all four operations should take $\Theta(\lg n)$ expected time.

% -----------------------------------------------------------------------------
\section{Implementation}

All algorithms were written from scratch in Java~21 without external
libraries. Nodes are objects with an \texttt{int} key and \texttt{left},
\texttt{right} and \texttt{p} references; order-statistic nodes add an
\texttt{int size}. The plain BST and the order-statistic tree are separate
classes so that the timings of one are not affected by the other's fields.

To keep timings meaningful, \WALK{} adds each key to a running sum instead of
printing it, since printing would cost far more than the walk. The tree height
is computed with a level-by-level traversal so that a chain of length $n$
cannot overflow the call stack.

% -----------------------------------------------------------------------------
\section{Experimental Methodology}

\subsection{Platform}

All runs were made on a cloud virtual machine with two vCPUs of an Intel Xeon
processor at \SI{2.80}{\giga\hertz}, \SI{32}{\kibi\byte} of L1 data cache and
\SI{1}{\mebi\byte} of L2 cache per core, \SI{33}{\mebi\byte} of shared L3
cache and \SI{8}{\giga\byte} of RAM, running Linux and OpenJDK~21.0.12 with a
\SI{4}{\gibi\byte} maximum heap. Each tree node occupies about
\SI{32}{\byte}, so a tree of $2^{15}$ nodes fills one core's L2 cache and a
tree of $2^{20}$ nodes fills the L3 cache.

\subsection{Input Generation}

Each tree was built from the keys $1, \dots, n$ in an order produced by a
Fisher--Yates shuffle, so every permutation is equally likely and every tree
is a randomly built BST in the sense of Section~12.4 of~\cite{clrs}. A fresh permutation was
drawn for every tree. Keys are distinct by construction.

Part~A used 37 values of $n$, namely $\operatorname{round}(2^{k/2})$ for
$k = 8, \dots, 44$, from $n = 16$ to $n = 2^{22} = 4\,194\,304$. These are
evenly spaced on a log scale and include every power of two in that range.
The number of trees per size was chosen so that every size received roughly
the same total work: 1000 trees for $n \le 2896$, falling to 10 trees from
$n = 2^{18}$ upwards, for 18\,612 trees in all. Each tree was built, its height
recorded, walked, and then destroyed.

Three supplementary runs isolated specific effects. The first built trees from
\emph{sorted} keys (chains) for $n$ from $2^6$ to $2^{15}$, five per size, to
time the walk on the worst-case shape. The second rebuilt random trees without
timing them and counted the nodes visited by building and by destroying
(30 trees per size up to $2^{16}$, five above). The third followed the height
of fifty trees of $2^{16}$ nodes at every halving of their size as they were
destroyed.

\subsection{Timing}

Times were taken with \texttt{System.nanoTime}. Node objects were allocated
before the timer started, because \INS{} takes an existing node as its
argument and allocation inside the timed loop also measured the garbage
collector copying the growing tree. Thirty untimed builds of 20\,000-node trees
preceded each run to let the JIT compiler finish optimising the code, and a
garbage collection was requested between sizes. A walk of a small tree takes
well under a microsecond, so the walk was repeated until about $2^{18}$ nodes
had been visited and the time divided by the number of repetitions; it was
timed on the first 30 trees of each size. The walk's key sum was accumulated
into a checksum printed at the end, so the compiler could not remove the walk
as dead code~\cite{georges}.

\subsection{Correctness}

Before any timing, a separate test built 200 random trees of sizes from 1 to
1000 with both tree types, then deleted half of the keys in random order and
finally destroyed each tree from the root. After every operation on trees of
up to 100 keys (and periodically on larger ones) it checked the BST property,
every parent pointer, and every \textit{size} field against a recount of the
subtree, and it checked \SEL{}$(i)$ and \RANK{} against a sorted list of the
keys for every $i$. All checks passed. To make sure the test could detect
errors, it was rerun on two deliberately broken versions of \DEL{}, one
omitting $y.\mathit{size} = z.\mathit{size}$ and one starting the decrement at
$z.p$ instead of $y.p$; both were caught with thousands of failures.

% -----------------------------------------------------------------------------
\section{Results: Binary Search Trees}

\subsection{Expected Height}

\begin{figure}[t]
\centering
\includegraphics[width=\columnwidth]{height.pdf}
\caption{Mean height of randomly built BSTs over 37 sizes, against the textbook
bound~\eqref{eq:clrs}, Reed's asymptotic expansion~\eqref{eq:reed} and the
minimum possible height. The shaded band is one standard deviation.}
\label{fig:height}
\end{figure}

\begin{table}[t]
\centering
\caption{Height of randomly built BSTs at selected sizes.}
\label{tab:height}
\begin{tabular}{rrrcrr}
\toprule
$n$ & trees & mean $h$ & range & s.d. & $h/\lg n$ \\
\midrule
$2^{4}$  & 1000 &  6.11 & 4--10  & 1.08 & 1.53 \\
$2^{8}$  & 1000 & 15.66 & 12--23 & 1.66 & 1.96 \\
$2^{12}$ &  814 & 26.63 & 22--34 & 2.04 & 2.22 \\
$2^{16}$ &   38 & 37.53 & 34--43 & 2.26 & 2.35 \\
$2^{20}$ &   10 & 49.60 & 47--52 & 2.01 & 2.48 \\
$2^{22}$ &   10 & 55.30 & 52--60 & 2.36 & 2.51 \\
\bottomrule
\end{tabular}
\end{table}

Fig.~\ref{fig:height} and Table~\ref{tab:height} show the measured heights.
Plotted against $\lg n$, the mean height is very close to a straight line,
which is exactly what $E[h] = O(\lg n)$ predicts. A least-squares fit over
$n \ge 2^{10}$ gives
\begin{equation}
\bar h \approx 2.856\lg n - 7.70,
\end{equation}
so each doubling of $n$ adds about 2.86 levels to the tree. This matches the
local slope of $2.87$ obtained from~\eqref{eq:reed} at the middle of the
range. The measured heights sit about five levels below Reed's expression,
a constant gap absorbed by its $O(1)$ term, and below the textbook
bound~\eqref{eq:clrs} at every size. The bound is therefore not just correct
but reasonably tight: at $n = 2^{22}$ it allows a mean height of 61.4,
against 55.3 measured.

The heights are also tightly concentrated. The standard deviation stayed
below 3.2 levels at every size, and no tree with more than 1000 keys had a
height above $3.2\lfloor\lg n\rfloor$, far from the $n-1$ of a chain. The ratio $h/\lg n$ rises
slowly with $n$ because of the negative constant in the fit; it should approach
$\alpha\ln 2 \approx 2.99$ only as $n \to \infty$.

\subsection{Building and Destroying}

\begin{figure*}[t]
\centering
\includegraphics[width=\textwidth]{build_destroy.pdf}
\caption{Building a random BST with $n$ calls to \INS{} and destroying it with
$n$ calls to \DEL{} on the root. (a) Mean running time. (b) Nodes visited,
measured in an untimed run, against the exact expectations~\eqref{eq:build}
and~\eqref{eq:destroy}; both are normalised by $n\lg n$. (c) Running time
divided by nodes visited.}
\label{fig:build}
\end{figure*}

Fig.~\ref{fig:build}(a) shows that building and destroying both grow slightly
faster than $n\lg n$ for small trees and noticeably faster for large ones.
Regression of $\log T$ on $\log n$ over $2^8 \le n \le 2^{22}$ gives exponents
of $1.42 \pm 0.04$ for building and $1.35 \pm 0.05$ for destroying, where an
exact $n\lg n$ curve would give 1.10 over the same range.

Fig.~\ref{fig:build}(b) shows this is not because the algorithms do more work
than expected. The number of nodes visited agrees with the exact
expectations~\eqref{eq:build} and~\eqref{eq:destroy} to within about 3\% at
every size from $2^8$ upward, and converges towards the predicted limits of
$1.386\,n\lg n$ for building and $0.693\,n\lg n$ for destroying (1.25 and 0.68
at $n = 2^{22}$, since the lower-order terms decay slowly). Destroying a tree
visits 0.53 to 0.62 times as many nodes as building it, matching the
prediction of one half.

The deletion trace supports the assumption behind
\eqref{eq:destroy}. Across fifty trees of $2^{16}$ nodes, the height at every
halving during destruction (Fig.~\ref{fig:deltrace}) stayed within 0.42 levels
of the mean height of a freshly built random tree of the same size, so
repeated root deletion does not unbalance the tree.

\begin{figure}[t]
\centering
\includegraphics[width=\columnwidth]{deltrace.pdf}
\caption{Height of a tree of $2^{16}$ keys as it is destroyed by deleting the
root, compared with freshly built random trees of the same size. Mean of 50
trees.}
\label{fig:deltrace}
\end{figure}

The extra growth in time comes entirely from the cost of each visit,
Fig.~\ref{fig:build}(c). Up to about $2^{15}$ nodes a visit costs a steady
6--9\,ns for building and 5--7\,ns for destroying. Beyond that the tree no
longer fits in the \SI{1}{\mebi\byte} L2 cache, and beyond $2^{20}$ nodes it no
longer fits in the L3 cache either; at $n = 2^{22}$ a visit costs 100\,ns when
building and 75\,ns when destroying. The path from the root to a random leaf
jumps between nodes allocated at unrelated times, so for a large tree almost
every step down is a cache miss. The algorithm's step count is
$\Theta(n\lg n)$ as theory says, but the price of a step is not constant.

\subsection{Inorder Tree Walk}

\begin{figure*}[t]
\centering
\includegraphics[width=\textwidth]{walk.pdf}
\caption{Time for one \WALK{}. Random BSTs with nodes allocated in insertion
order (the normal case), random BSTs with the same shape but nodes allocated
in key order, and chains built from sorted keys. (a) Mean time per walk.
(b) Time per node.}
\label{fig:walk}
\end{figure*}

\WALK{} visits each node exactly once, so its operation count is exactly $n$.
On chains the measured time grew with exponent $1.04 \pm 0.08$
(Fig.~\ref{fig:walk}), confirming $\Theta(n)$ even on the worst-case shape.
Building those chains, by contrast, grew with exponent $1.94 \pm 0.08$: sorted
input makes every \INS{} walk the whole chain, so building costs
$\Theta(n^2)$ and took 1.35\,s for only $2^{15}$ keys, compared with
\SI{5}{\milli\second} for a random tree of that size.

On random trees the walk scaled with exponent $1.41 \pm 0.03$, and the time
per node rose from 4.4\,ns at $n = 2^8$ to 135\,ns at $n = 2^{22}$. The walk
visits nodes in key order, but they sit in memory in insertion order, which
is random; once the tree exceeds the caches, every visit is a cache miss.
To confirm this, the experiment was repeated with the nodes allocated in key
order but inserted in the same random order, giving trees of the same random
shape whose nodes lie in memory in the order the walk visits them. These
walks scaled with exponent $1.23 \pm 0.02$ and were $5.3\times$ faster at
$n = 2^{22}$ (25.7\,ns per node). They do not reach the chain's $n^{1.04}$
because the recursion still moves up and down a tree of height about 55 and
the JVM may relocate objects during garbage collection, so the layout is only
approximately sequential. The walk performs the same $n$ visits in every
case; only the memory layout changed.

% -----------------------------------------------------------------------------
\section{Results: Order-Statistic Trees}

\subsection{Setup}

For part B, both a plain BST and an order-statistic tree were built from the
same $n$ shuffled keys, for 25 sizes from $n = 2^{10}$ to $2^{22}$ (40 trials
per size up to $2^{16}$, 20 up to $2^{19}$, and 10 above). In each trial
$m = 1000$ further keys were then inserted into each tree, 1000
\SEL{}$(T.\mathit{root}, i)$ calls were made with $i$ uniform over the current
size, 1000 \RANK{} calls were made on uniformly chosen nodes, and 1000
uniformly chosen nodes were deleted from each tree. Finally both trees were
destroyed from the root. Each batch of 1000 operations was timed as a whole
and divided by 1000. The two tree types alternated in going first, and both
were given the same memory layout (nodes allocated in insertion order, with a
garbage collection before timing), so that differences between them come from
the algorithms and not from where their nodes happen to sit in memory.

\subsection{Order-Statistic Operations}

\begin{figure*}[t]
\centering
\includegraphics[width=\textwidth]{os_ops.pdf}
\caption{The four operations on order-statistic trees. (a) Mean time per
operation, with one standard deviation. (b) Time divided by the expected
number of nodes on a root-to-node path~\eqref{eq:depth}.}
\label{fig:osops}
\end{figure*}

\begin{table}[t]
\centering
\caption{Mean time per operation (ns) on plain BSTs and order-statistic (OS)
trees.}
\label{tab:ops}
\setlength{\tabcolsep}{4.5pt}
\begin{tabular}{rrrrrrr}
\toprule
 & \multicolumn{2}{c}{\INS{}} & \multicolumn{2}{c}{\DEL{}} & & \\
\cmidrule(lr){2-3}\cmidrule(lr){4-5}
$n$ & BST & OS & BST & OS & \SEL{} & \RANK{} \\
\midrule
$2^{10}$ &  106 &  128 &  116 &  165 &  354 &  180 \\
$2^{13}$ &  210 &  274 &  149 &  207 &  208 &  254 \\
$2^{16}$ &  904 &  779 &  311 &  596 &  664 &  444 \\
$2^{19}$ & 1894 & 1927 & 1008 & 2078 & 1765 & 1184 \\
$2^{22}$ & 3324 & 3091 & 1180 & 2796 & 3195 & 2243 \\
\bottomrule
\end{tabular}
\end{table}

Fig.~\ref{fig:osops}(a) and Table~\ref{tab:ops} show the mean cost of each
operation. Between $n = 2^{10}$ and $n = 2^{22}$ the expected root-to-node
path~\eqref{eq:depth} grows from 13.4 to 28.7 nodes, a factor of 2.1, while
the measured times grew by factors of 9 to 24. Dividing each time by the
expected path length (Fig.~\ref{fig:osops}(b)) explains the difference. Up to
about $2^{13}$ keys every operation costs a roughly constant 10--17\,ns per
node on its path. From there the cost per node climbs to 80--110\,ns at
$n = 2^{22}$, the same rise seen in the cost per visit when building and
destroying trees (Fig.~\ref{fig:build}(c)). The four operations therefore do
$\Theta(\lg n)$ work each, as expected, and the extra growth is again the
memory system. Individual batches are short (0.1--0.3\,ms for small trees),
so the small sizes are noisy, for example \SEL{} at $n = 2^{10}$; the error
bars show the spread.

\RANK{} was consistently the cheapest of the four, about 30\% below the others
per node at large $n$. A likely reason is that \SEL{} must read the left
child's size before it knows which way to go, so each step waits on a memory
access, whereas the next node in \RANK{} is simply the parent, and the sibling
size it reads does not decide the path. The processor can therefore overlap
more of \RANK{}'s memory accesses.

\subsection{Cost of Maintaining the Size Attribute}

\begin{figure*}[t]
\centering
\includegraphics[width=\textwidth]{os_vs_bst.pdf}
\caption{Plain BST against order-statistic tree on identical keys.
(a) Mean time per \INS{} and per \DEL{} of a random node. (b) Ratio of
order-statistic tree time to plain BST time, for single operations and for
whole builds and destroys.}
\label{fig:osbst}
\end{figure*}

Fig.~\ref{fig:osbst} compares the two trees directly. For insertion the extra
work is invisible: the order-statistic \INS{} took on average 1.02 times as
long as the plain one, and 1.00 times for $n > 2^{16}$. The size increment
happens on a node the search has just loaded to compare keys, so it costs
almost nothing. Whole builds behaved the same way (ratio 1.09 overall, 1.02
above $2^{16}$), as did destroying from the root (1.13 overall, 1.06 above
$2^{16}$). In the latter case the decrement walk starts at the successor's
parent and climbs the left spine that \textsc{Tree-Minimum} has just walked
down, so those nodes are still in the cache.

Deleting a random node is different. The order-statistic \DEL{} took 1.85 times
as long as the plain one on average, rising from about 1.6 for
$n \le 2^{16}$ to 2.1 above and reaching 2.45. The untimed run explains why.
Given a pointer to $z$, plain \DEL{} only touches $z$ and, if $z$ has two
children, the path to its successor. Over every node of the random trees, this
averaged 1.75 nodes at every size from $2^{10}$ to $2^{22}$, so on a random
tree plain deletion of a known node takes expected constant time. The
augmented version must also visit every ancestor of the vacated position,
which averages 12.0 nodes at $n = 2^{10}$ and 28.6 at $n = 2^{22}$, making it
$\Theta(\lg n)$. The measured ratio is much smaller than the ratio of node
counts because the nodes near the root lie on every path and stay in the
cache, so most of the upward walk is cheap; the cost is dominated by the few
misses near the bottom of the tree, which both versions pay. The plain
\DEL{} still slowed from 116\,ns to 1180\,ns over the range for the same reason:
the one or two nodes it touches are scattered in memory.

% -----------------------------------------------------------------------------
\section{Discussion}

Three conclusions follow from these measurements.

First, the theory was confirmed wherever it makes a prediction about work done.
The expected height of a randomly built BST grew linearly in $\lg n$ with the
slope predicted by~\eqref{eq:reed}, inside the bound of Theorem~12.4, and with
little variation between trees. The numbers of nodes visited by building and
destroying matched their exact expectations to within about 3\%, the tree stayed
random as it was destroyed, \WALK{} performed $n$ visits on every shape, and
the order-statistic operations followed paths of the expected
$\Theta(\lg n)$ length.

Second, running time did not follow the work as closely as in a flat-memory
model. The cost of one step down a tree rose about fifteenfold between small
trees and trees of four million nodes, because a pointer-based tree with a
random shape spreads consecutive steps across unrelated parts of memory.
This made $\Theta(n\lg n)$ procedures look like $n^{1.4}$ and a $\Theta(n)$
walk look like $n^{1.41}$. Changing only where the nodes were allocated gave
the walk a $5.3\times$ speedup. Separating operation counts from times was
what made it possible to tell these two effects apart.

Third, the size augmentation is cheap exactly where the original operation
already visits the affected nodes. Insertion and \SEL{} walk the root path
anyway, so they gained little or nothing. Deletion of a known node does not,
so maintaining sizes turned an expected constant-time operation into a
$\Theta(\lg n)$ one and doubled its time on large trees. The textbook's
red-black version pays for the same upward walk, but there it sits inside a
deletion that is already $O(\lg n)$ in the worst case.

Some limitations apply. All results come from one virtual machine, so the
constants and the sizes where the cache effects begin will differ on other
hardware, although the operation counts and the height results do not depend
on the machine. The JVM moves objects during garbage collection, so memory
layout was controlled only approximately. Finally, every tree was built from
random insertions only. Long sequences that mix insertions and deletions are
known to make BSTs less balanced than a random tree~\cite{eppinger}, so the
height results here apply to randomly built trees and not to trees with an
arbitrary history of updates.

% -----------------------------------------------------------------------------
\section{Conclusion}

\INS{}, \DEL{} and \WALK{} were implemented as in~\cite{clrs}, along with an
order-statistic tree on a plain BST supporting \SEL{} and \RANK{}. Over
18\,612 randomly built trees of up to $2^{22}$ keys, the mean height grew as
$2.86\lg n - 7.7$, confirming the $O(\lg n)$ expected height of Theorem~12.4
and staying below its bound at every size. Building and destroying visited
$\Theta(n\lg n)$ nodes in close agreement with their exact expectations, and
the inorder walk visited exactly $n$. Size maintenance in \INS{} costs one
increment per node on the search path and in \DEL{} a walk from the vacated
position to the root; this left insertion as fast as before but made deletion
about twice as slow on large trees. Across all experiments the operation
counts followed the theory closely, while the running times grew faster
because the cost of each step rose as the trees outgrew the processor caches.

\begin{thebibliography}{9}

\bibitem{clrs}
T.~H. Cormen, C.~E. Leiserson, R.~L. Rivest, and C.~Stein, \textit{Introduction
to Algorithms}, 3rd~ed. Cambridge, MA, USA: MIT Press, 2009.

\bibitem{devroye}
L.~Devroye, ``A note on the height of binary search trees,''
\textit{J. ACM}, vol.~33, no.~3, pp.~489--498, 1986.

\bibitem{reed}
B.~Reed, ``The height of a random binary search tree,'' \textit{J. ACM},
vol.~50, no.~3, pp.~306--332, 2003.

\bibitem{knuth}
D.~E. Knuth, \textit{The Art of Computer Programming, Vol.~3: Sorting and
Searching}, 2nd~ed. Reading, MA, USA: Addison-Wesley, 1998.

\bibitem{hibbard}
T.~N. Hibbard, ``Some combinatorial properties of certain trees with
applications to searching and sorting,'' \textit{J. ACM}, vol.~9, no.~1,
pp.~13--28, 1962.

\bibitem{eppinger}
J.~L. Eppinger, ``An empirical study of insertion and deletion in binary
search trees,'' \textit{Commun. ACM}, vol.~26, no.~9, pp.~663--669, 1983.

\bibitem{georges}
A.~Georges, D.~Buytaert, and L.~Eeckhout, ``Statistically rigorous Java
performance evaluation,'' in \textit{Proc. OOPSLA}, 2007, pp.~57--76.

\end{thebibliography}

\end{document}
